%% kapa: exam=1 round=29 year=2018 period=2 form=A subject=appraisal_law subject_ko=감정평가관계법규 seq=36 number=36
%% source: Q-Net 감정평가사 기출문제 (cst00309) artlSeq=5202808 file=extracted/1008/5202808/2018년 제29회 감정평가사 1차시험 기출문제/2018년 제29회 감정평가사 1차시험 기출문제/2교시 A형.pdf
%% conversion: kordoc 4.24.51 -> tools/gen1.py (auto); review=proofread-sonnet; flags=none
\begin{question}{36}
동산ㆍ채권 등의 담보에 관한 법률상 담보등기에 관한 설명으로 옳지 않은 것은?
\begin{choices}
  \item 담보등기는 법률에 다른 규정이 없으면 등기권리자와 등기의무자가 공동으로 신청한다.
  \item 등기관이 등기를 마친 경우 그 등기는 접수한 때부터 효력을 발생한다.
  \item 「동산ㆍ채권 등의 담보에 관한 법률」에 따른 담보권의 존속기간은 5년을 초과할 수 없으나, 5년을 초과하지 않는 기간으로 이를 갱신할 수 있다.
  \item 등기관의 결정 또는 처분에 대한 이의신청은 집행정지의 효력이 없다.
  \item 등기관의 결정 또는 처분에 이의가 있는 자는 새로운 사실이나 새로운 증거방법을 근거로 관할 지방법원에 이의신청을 할 수 있다.
\end{choices}
\end{question}
\answer{5}
